%; whizzy chapter -dvi
% -initex iniptex -latex platex -format platex -bibtex jbibtex -fmt fmt
% 以上 whizzytex を使用する場合の設定。
 
%     Tokyo Debian Meeting resources
%     Copyright (C) 2012 Junichi Uekawa
%     Copyright (C) 2011, 2015 Nobuhiro Iwamatsu

%     This program is free software; you can redistribute it and/or modify
%     it under the terms of the GNU General Public License as published by
%     the Free Software Foundation; either version 2 of the License, or
%     (at your option) any later version.

%     This program is distributed in the hope that it will be useful,
%     but WITHOUT ANY WARRANTY; without even the implied warranty of
%     MERCHANTABILITY or FITNESS FOR A PARTICULAR PURPOSE.  See the
%     GNU General Public License for more details.

%     You should have received a copy of the GNU General Public License
%     along with this program; if not, write to the Free Software
%     Foundation, Inc., 51 Franklin St, Fifth Floor, Boston, MA  02110-1301 USA

%  preview (shell-command (concat "evince " (replace-regexp-in-string "tex$" "pdf"(buffer-file-name)) "&"))

%%ここからヘッダ開始。

\documentclass[mingoth,a4paper]{jsarticle}
\usepackage{monthlyreport}
% 日付を定義する、毎月変わります。
\newcommand{\debmtgyear}{2017}
\newcommand{\debmtgmonth}{11}
\newcommand{\debmtgdate}{18}
% started from zero:
% (let ((year 2013) (month 7)) (+ (* (- year 2005) 12) month -1))
\newcommand{\debmtgnumber}{157}

% Needed to import pandoc-generated LaTeX documents.
% See https://stackoverflow.com/questions/40438037/tightlist-error-using-pandoc-with-markdown
\providecommand{\tightlist}{%
  \setlength{\itemsep}{0pt}\setlength{\parskip}{0pt}}

% tikz picture の為のマクロ設定
\usepackage[dvipdfmx]{graphicx}
\usepackage{tikz}

\begin{document}

\begin{titlepage}
\thispagestyle{empty}
% タイトルページ:編集必要な部分は最初のマクロに飛ばすこと

\vspace*{-2cm}
第\debmtgnumber{}回 東京エリア Debian 勉強会資料\\
\hspace*{-2cm}
\includegraphics{image2012-natsu/dotdeb.pdf}\\
\hfill{}\debmtgyear{}年\debmtgmonth{}月\debmtgdate{}日

% ここはアップデートすること
% 全角文字にしないとフォントのサイズが合わないので注意
\rotatebox{10}{\fontsize{30}{30} {\gt DPL来日、リリース再考 }}\\

\vspace*{-2cm}
\hfill{}\includegraphics[height=6cm]{image200502/openlogo-nd.eps}
\end{titlepage}

\newpage

\begin{minipage}[b]{0.2\hsize}
 \definecolor{titleback}{gray}{0.9}
 \colorbox{titleback}{\rotatebox{90}{\fontsize{80}{80} {\gt デビアン勉強会} }}
\end{minipage}
\begin{minipage}[b]{0.8\hsize}
\hrule
\vspace{2mm}
\hrule
\begin{multicols}{2}
\tableofcontents
\end{multicols}
\vspace{2mm}
\hrule
\end{minipage}

\dancersection{最近のDebian関連のミーティング報告}{杉本 典充}

\subsection{第156回東京エリアDebian勉強会}

2017年10月21日(土)に第156回東京エリアDebian勉強会を開催しました。会場は東銀座にある朝日ネットさんをお借りして行いました。参加者は5名でした（1名は飛び入り参加）。

発表は吉野さんによる「caffとmail-transport-agent」の1本でした。

「caffとmail-transport-agent」では、caffの紹介とメール送信するMTAは何がおすすめか紹介する発表でした。メール送信はSMTPの25/TCP ポートへ接続する方法とsendmailコマンドを起動する方法があり、その特徴を整理しました。またdebianで提供する仮想パッケージ「mail-transport-agent」の説明を行い、mail-transport-agentの中でおすすめはgmail等の外部メールサービスを使ってメール送信できるmsmtp{\footnote{\url{https://wiki.debian.org/msmtp}}}であるとの紹介がありました。

その後ハックタイムを行い、作業内容を共有しました。


\dancersection{事前課題}{杉本 典充}

今回の事前課題は以下です:
\begin{enumerate}
\item Hack Timeは何をしますか
\item DPLへぜひとも質問をお願いします（複数記述可）
\end{enumerate}
この課題に対して提出いただいた内容は以下です。
\begin{multicols}{2}
{\small
\begin{prework}{ ftake }
  \begin{enumerate}
  \item HackMD $B$r%$%s%9%H!<%k$7$^$9(B
  \item Which do you prefer community management or development? How long can you use time for community activity in a day or a week?
  \end{enumerate}
\end{prework}

\begin{prework}{ yoshitake }
  \begin{enumerate}
  \item $B%;%C%7%g%s$N$_$G$9$,!";22C$5$;$F$$$?$@$-$^$7$?$i9,$$$G$9!#(B
  \item Good to see you again ;)
  \end{enumerate}
\end{prework}

\begin{prework}{ khibino }
  \begin{enumerate}
  \item cabal-debian $B$N(B proposed-update $B$X$N(B upload $B=`Hw(B
  \item (thinking...)
  \end{enumerate}
\end{prework}

\begin{prework}{ ghanshyamman }
  \begin{enumerate}
  \item Open source developer
  \item How is the Debian release model experience. Any feedback or key thing to share for other software.
  \end{enumerate}
\end{prework}

\begin{prework}{ yy\_y\_ja\_jp }
  \begin{enumerate}
  \item $B%Q%C%1!<%899?7!$(BDDTSS
  \item Outreach -- are there any ideas for women and young people to attract Debian?
  \item init systems -- I use sysvinit instead of systemd. I feel some DDs have some intention to make developers/users uncomfortable who support other init systems than systemd. In my opinion it's not a good idea for developers' or users' diversity and universality. Around the GR period some DDs have left Debian. I think it's unfortunate. What do you think about this?
  \end{enumerate}
\end{prework}

\begin{prework}{ Roger Shimizu }
  \begin{enumerate}
  \item pkg maintainence of torbrowser-launcher
  \item I hope more infrastructure, such as buildd, ftp-master server can move to asia, or even japan! But how can local community help?
  \end{enumerate}
\end{prework}

\begin{prework}{ dictoss }
  \begin{enumerate}
  \item gcc$B$N(Bpie$B%*%W%7%g%s$rD4$Y$FH/I==`Hw$r$9$k(B
  \item Please tell me popular sub project of debian. I dot't find sub project guide page in debian wiki.
  \end{enumerate}
\end{prework}

\begin{prework}{ koedoyoshida }
  \begin{enumerate}
  \item $BL$Dj(B
  \item (thinking...)
  \end{enumerate}
\end{prework}

\begin{prework}{ ysaito }
  \begin{enumerate}
  \item $B%I%-%e%a%s%HFI$_(B
  \item The DPL may largely be entangled in daily business to keep the project going. Do you have a personal dream or vision on what Debian should be like in a distant future?
  \end{enumerate}
\end{prework}

\begin{prework}{ masayukig }
  \begin{enumerate}
  \item Read the docs and setup the Debian openstack cloud image to my own ``cloud''
  \item What are challenges of the Debian community?
  \item Do you have any expectations for Japanese user group?
  \item What technologies are you interested in, recently?
  \end{enumerate}
\end{prework}

\begin{prework}{ BurnDuck (nabaua) }
  \begin{enumerate}
  \item tensorflow
  \item how to update debian 9
  \end{enumerate}
\end{prework}

\begin{prework}{ $B$O$7$b$H$^$5(B }
  \begin{enumerate}
  \item GPU$BF~$j%N!<%H(BPC$B$K(BDebian$B$O%$%s%9%H!<%k$G$-$k$+!)!J$r$d$C$F$_$k$+$b!K(B
  \item How about in Yokohama ? ($B2#IM$N46A[$r$<$R!#(B)
  \end{enumerate}
\end{prework}

\begin{prework}{ NOKUBI Takatsugu }
  \begin{enumerate}
  \item $B<!$NJY6/2q$N;qNA:n@.(B
  \item Do you know the age bracket of Debian Developeres and Maintaineres ? $B!J:#3+H/<T$NG/NpAX$O$I$&$J$C$F$$$k$+!#!K(B
  \item Are young Debian Developeres and Maintaineres increasing ? $B!J<c$$?M$OA}$($F$$$k$N$+$I$&$+!K(B
  \end{enumerate}
\end{prework}

\begin{prework}{ henrich }
  \begin{enumerate}
  \item $B%G%#%9%+%C%7%g%s(B
  \item What would you like to achieve in your DPL term?
  \end{enumerate}
\end{prework}

}
\end{multicols}

%\dancersection{Debian Trivia Quiz}{杉本 典充}
%
%Debianの昨今の話題についてのQuizです。
%
%今回の出題範囲は\url{debian-devel-announce@lists.debian.org} や \url{debian-news@lists.debian.org}などに投稿された内容からです。
%
%\begin{multicols}{2}
%%; whizzy-master ../debianmeetingresume201311.tex
% $B0J>e$N@_Dj$r$7$F$$$k$?$a!"$3$N%U%!%$%k$G(B M-x whizzytex $B$9$k$H!"(Bwhizzytex$B$,MxMQ$G$-$^$9!#(B
%

\santaku
{quiz}
{A}
{B}
{C}
{C}
{desc}

%\end{multicols}


% % (query-replace-regexp "<.*?>" "")
% % (query-replace-regexp "^[	 ]\+" "")

%-------------------------------------------------------------------------------
\dancersection{Debianのリリースを再考する}{Yamane Hideki}
%-------------------------------------------------------------------------------

発表資料を参照ください。

%-------------------------------------------------------------------------------
\dancersection{DPLへの質問コーナー}{Chris Lamb}
%-------------------------------------------------------------------------------

事前課題を参照ください。

%
% 冊子にするために、4の倍数にする必要がある。
% そのための調整
%\dancersection{メモ}{}
\mbox{}\newpage

\vspace*{15cm}
\hrule
\vspace{2mm}
\includegraphics[width=2cm]{image200502/openlogo-nd.eps}
\noindent \Large \bf Debian 勉強会資料\\
\noindent \normalfont \debmtgyear{}年\debmtgmonth{}月\debmtgdate{}日 \hspace{5mm}  初版第1刷発行\\
\noindent \normalfont 東京エリア Debian 勉強会 （編集・印刷・発行）\\
\hrule

\end{document}
